\subsubsection{Ramanujan Symmetry and Horizon Compactification: Positivity Certificates, Spectral Bounds, and Information Ledger Framework for Riemann Hypothesis Zeros}\label{subsubsec:xi-ramanujan-horizon-ledger-framework}
This subsection provides a synoptic reduction, in the format of an ``auditable certificate chain,'' for the \(\Xi\)/unitary slice audit content of this section: its goal is not to introduce new assumptions or new proof pivots, but rather to rearrange the already closed numbered conclusions by their interface relations into a seamless closed loop, so that the three-layer constraints of ``positivity--spectral bounds--information budget'' are aligned within a single completion--horizon--certificate semantics.
Within this framework, recursive addressing should be understood as an \emph{endogenization mechanism for readable prefixes}: when the address object is not given, the evaluation function \(\ind_{C_a}\) fails at the object/type level, the read partial function is undefined and denoted \(\mathsf{NULL}\) to indicate non-addressability; once the address is given, the recursively generated cylinder event family does not expand the existing visible \(\sigma\)-algebra, but merely reselects and reorganizes visible events within the old-layer visible domain, whence \(\mathcal{G}^{(L+1)}\subseteq\mathcal{G}^{(L)}\) (Proposition \ref{prop:null-as-local-section-obstruction}, \ref{prop:recursive-addressing-refinement}, and Corollary \ref{cor:recursive-addressing-visible-sigma-nonexpansion}).
Therefore, the ``certificates/falsifiers/lower bounds'' stated below should all be understood as: given the address (prefix) and the visible domain constraint, the output must be encoded as a finite object and be offline-verifiable by a third party; any data channel that has not been explicitly addressed and entered into the certificate object must not affect the verification conclusion at the semantic level.

\paragraph{Three-Step Closed Chain of the Diskified Horizon Interface: Unitary Slice, Endpoint \texorpdfstring{$-1$}{-1}, and Inner Factor Hair}
Via the Cayley compactification
\(
w=\mathcal{C}(t)=\frac{1+\mathrm{i}t}{1-\mathrm{i}t}
\),
the two degrees of freedom of \(t=\gamma+\mathrm{i}\delta\in\mathbb{H}\) are transported to the unit disk \(\mathbb{D}\); then off-critical-strip zeros
\(\rho=\beta+\mathrm{i}\gamma\) (taking \(\beta<\tfrac12\), \(\delta:=\tfrac12-\beta>0\))
are strictly equivalent to disk-interior zero events, and the exponentialized Busemann/Poisson coordinate at the endpoint \(-1\),
\(\mathcal{B}_{-1}(w)=\frac{1-|w|^{2}}{|1+w|^{2}}\), realizes a closed-form readout of the offset:
\(
\delta=\mathcal{B}_{-1}(w_\rho)
\)
and
\(
1-|w_\rho|^{2}=\delta|1+w_\rho|^{2}
\)
(Theorem \ref{thm:app-rh-iff-disk-zero-free}, Proposition \ref{prop:app-busemann-poisson-minus1}, Corollary \ref{cor:app-offcritical-horocycle-busemann}).
Within this closed context, the structural main chain needed in this paper can be compressed into three interface statements (subsequent paragraphs merely provide numbered expansions and quantification):
{\renewcommand{\labelenumi}{(\roman{enumi})}
\begin{enumerate}
  \item \textbf{Ramanujan-type spectral symmetry (phase closure/unitary slice)} fixes the visible readout on the unitary slice \(|u|=1\) and constrains verifiable zero events to this fixed slice; off-slice information can only enter the certificate domain in the form of defects/falsifiers (e.g., Theorem \ref{thm:phase-lift-fixed-slice} and \S\ref{subsubsec:selfdual-blaschke-defect-fixed-slice}).
  \item \textbf{Endpoint \texorpdfstring{$-1$}{-1} horizon geometry} compresses the off-critical-line offset \(\delta\) into the horocycle depth coordinate \(\mathcal{B}_{-1}\) and forces it to enter the endpoint thin layer via a second-order radial channel (Corollary \ref{cor:app-offcritical-horocycle-busemann}; see also Theorem \ref{thm:xi-endpoint-compression-precision}).
  \item \textbf{Inner factor carrier (Blaschke/singular inner factor)} provides the unique carrier for ``horizon hair'': disk-interior zeros and singular support can only be carried by the inner factor \(BS\), while the outer factor is determined solely by the boundary modulus; under Hardy/Smirnov conditions, RH is equivalent to the vanishing of disk-interior zero data (\(B\equiv 1\)), and under the no-singular-inner-factor assumption, it is further equivalent to the endpoint phase current certificate being zero (Appendix \S\ref{subsec:unit-disk-inner-outer}).
\end{enumerate}}
From this, two types of auditable corollaries follow immediately: on the one hand, RH is equivalent to Carath\'eodory positivity of the horizon logarithmic derivative and equivalent to the Toeplitz--PSD hierarchy (Appendix Theorem \ref{thm:app-horizon-toeplitz-lmi}), so any deviation must yield a finite-dimensional falsifier at some finite level; on the other hand, the endpoint second-order radial compression and Toeplitz detection threshold yield infeasibility bounds for finite-depth certificates (Theorem \ref{thm:xi-toeplitz-length-vs-boundary-depth}).
In the semantics of this paper, the closed chain ``finite-dimensional falsifier \(\Rightarrow\) resonance reconstruction \(\Rightarrow\) horocycle depth estimate \(\Rightarrow\) strictification invariant decomposition'' has been formalized by Theorem \ref{thm:xi-toeplitz-failure-closed-chain}.

\paragraph{Bulk--Horizon Compression/Certificate Framework: Interface Unification of Five Themes (Expository; Not Used as Inference Premises)}
The above interface chain can also be understood at the expository level as two-sided readouts of the same type of object: when the completed object is compressed into the unit disk under the self-dual constraint via Cayley compactification (Appendix \ref{app:unit-circle-phase-gate}), its core difficulty can be rigorously stated as a \textbf{horizon positivity problem for self-dual disk-completed objects}.
Here ``horizon'' refers to the boundary readout on the fixed slice \(|w|=1\) (unitary slice), and ``positivity'' means the horizon logarithmic derivative function belongs to the Carath\'eodory class and therefore possesses a Herglotz representation: there exists a finite positive measure \(\mu\) on the unit circle (denoted \(\mu_{\zeta}\) in the \(\zeta\) case), whose Fourier moments
\(
c_n=\int_{\TT}\xi^{-n}\,d\mu(\xi)
\)
yield a sequence of auditable horizon multipole moments (i.e., ``hair coefficients''), generating the Toeplitz--PSD hierarchy
\(
\mathbf{T}_N=(c_{j-k})_{0\le j,k\le N}\succeq 0
\)
(Theorem \ref{thm:app-horizon-toeplitz-lmi}; see also Appendix \ref{subsubsec:app-horizon-weyl-carath-toeplitz} for the measure/moment interface).
Under the same disk completion interface, ``disk-interior zero exclusion'' can also be equivalently compressed into the horizon purity criterion via the Jensen defect \(D_{\zeta}(\varrho)\) (Theorem \ref{thm:xi-horizon-purity-criterion}).
Therefore, the core of the unification is not analogy but rather the spectral structure determined by the positivity/positive semidefinite hierarchy that is level-by-level offline-verifiable on the unit circle boundary: failure of the hierarchy at any finite level provides a finite falsifier, while its full validity is equivalent to the ``horizon hair measure'' being positive.

Under this unified object, the five themes that recur throughout this paper---Riemann zeros, Ramanujan-type symmetry, black hole horizons, fractals, and \(d=6\) hypercubes (golden scale)---can be understood as different cross-sections of the same positivity problem (none introducing new inference premises):
\begin{itemize}
  \item \textbf{Riemann zeros:} RH \(\Leftrightarrow\) horizon positivity \(\Leftrightarrow\) Toeplitz--PSD hierarchy (Theorem \ref{thm:app-horizon-toeplitz-lmi}).
  \item \textbf{Ramanujan-type symmetry:} Phase closure fixes the visible domain to the unitary slice and indexes off-slice point spectrum defects as the Blaschke defect \(\mathfrak d_{\mathrm{Bl}}\); ``no hair'' is here equivalent to \(\mathfrak d_{\mathrm{Bl}}=0\) (Theorem \ref{thm:blaschke-defect-zero-fixed-slice-support} and Theorem \ref{thm:blaschke-point-spectrum-correspondence}).
  \item \textbf{Black hole horizon:} The endpoint \(w=-1\) corresponds to the compactification point \(t\to\infty\); the off-slice offset enters the horizon depth \(h(\rho)=1-|w_\rho|^2\) via a quadratic radial channel (Theorem \ref{thm:xi-offcritical-quadratic-radial-compression}).
  \item \textbf{Fractals:} The R\'enyi endpoint constant \(D_\infty\) of the folding pushforward measure provides the power-law scale for defect visibility (Proposition \ref{prop:pom-rq-qinfty-endpoint}; Corollary \ref{cor:xi-bulk-horizon-defect-visibility-Dinfty}), thereby furnishing an auditable fingerprint for the microstructure of the horizon measure.
  \item \textbf{\(d=6\) hypercube (golden scale):} The bulk side is organized by the \(d=6\) multi-channel potential and the reachable mean convex body (rotation/mean set) (\S\ref{subsec:pom-potential-convex-geometry}), while the window-\(6\) \(21\)-mode root--Cartan dictionary and Weyl orbit compression provide a finite-dimensional coarse-graining basis for the horizon hair moment data (Theorem \ref{thm:window6-unified-skeleton}, Theorem \ref{thm:boundary-m10-b3c3-rigid-four-block-split}, and Corollary \ref{cor:fold6-weyl-two-orbit-compression}).
\end{itemize}

\begin{proposition}[Window-\(6\) \(21\)-Mode Coarse-Graining and Weyl Compression: Low-Dimensional Parametrization of the Horizon Toeplitz Hierarchy]\label{prop:xi-horizon-21mode-weyl-compression}
At the boundary uplift layer of the window-\(6\) anchor, the \(B_3/C_3\) root--Cartan dictionary yields the rigid modal decomposition
\(
21=18\oplus 3
\):
\(18\) root directions and \(3\) Cartan directions (Theorem \ref{thm:boundary-m10-b3c3-rigid-four-block-split}; see also conclusion group (1) of Theorem \ref{thm:window6-unified-skeleton}).
If only Weyl-invariant statistical compression is permitted, the weights of the \(18\) root directions are at most a \(2\)-dimensional parameter family under coarse splitting (short/long two orbits) (Corollary \ref{cor:fold6-weyl-two-orbit-compression}), and at most a \(3\)-dimensional parameter family under fine splitting
\(
18=6\ (\mathrm{short})\oplus 6\ (A_2(+))\oplus 6\ (A_2(-))
\)
(the rigid block splitting of Theorem \ref{thm:boundary-m10-b3c3-rigid-four-block-split}).
Therefore, together with the \(3\) Cartan weights, the Weyl-invariant \(21\)-mode statistical family is at most a \(5\)- (coarse splitting) or \(6\)- (fine splitting) dimensional parameter family.
\par
Furthermore, if a horizon multipole moment sequence \(\{c_n\}_{n\in\ZZ}\) is given by the Fourier moments of a finite positive measure \(\mu\) on the unit circle,
\(
c_n=\int_{\TT}\xi^{-n}\,d\mu(\xi)
\),
and \(\mu\) is linearly parametrized on the above \(21\)-mode skeleton with these weights (e.g., as a finite approximation of the cut-and-project identification in \S\ref{subsubsec:discussion-bulk-horizon-dictionary}; see Conjecture \ref{conj:discussion-horizon-measure-cutproject}),
then for any \(N\ge 0\), the Toeplitz principal submatrix
\(
\mathbf{T}_N=(c_{j-k})_{0\le j,k\le N}
\)
is an affine function of these parameters; hence the constraint \(\mathbf{T}_N\succeq 0\) at each level is a linear matrix inequality (LMI) in at most \(6\) parameters, and equivalently provides a standard semidefinite feasibility (SDP) interface.
\end{proposition}
\begin{proof}
The Weyl orbit splitting and rigid block structure on root directions are given by Theorem \ref{thm:boundary-m10-b3c3-rigid-four-block-split}; in particular, Weyl-invariant weights must be constant on each orbit (or each rigid block), so the root direction weight degrees of freedom are compressed to \(2\) (short/long) or \(3\) (short/\(A_2(+)\)/\(A_2(-)\)) scalars, and adding the \(3\) Cartan weights yields the total parameter dimension upper bound.
\par
On the other hand, the Fourier moment \(c_n=\int_{\TT}\xi^{-n}\,d\mu(\xi)\) is linear in the measure \(\mu\); if \(\mu\) is linearly generated by finite modal weights, then each \(c_n\) and the Toeplitz matrix \(\mathbf{T}_N\) formed from them are affine functions of these weights. Hence \(\mathbf{T}_N\succeq 0\) is an LMI in the parameters, corresponding to a standard SDP feasibility constraint. \qed
\end{proof}

At the expository analogy level, the unit circle boundary can be understood as a ``black hole horizon''-style boundary slice: the second-order radial compression at the endpoint layer pushes high-height offsets into an extremely thin boundary layer, and the induced macrostate-scale quadratic scaling (Theorem \ref{thm:xi-bulk-horizon-phi-area-law}) corresponds to a strict ``area-type'' cost law.
This paper uses only the analytic closed forms of this compression and the certificate cost corollaries, and does not use the analogy as a proof premise.
The following three main results compress this expository picture into directly citable quantitative propositions; their proofs depend only on the numbered conclusions already given in this paper.

\paragraph{Three Main Results: \texorpdfstring{$\varphi$}{phi}-Area Law, Endpoint Fractal Index, and Toeplitz Certificate Transfer}
\begin{theorem}[\texorpdfstring{$\varphi$}{phi}-Area Law: Minimal Visible Resolution of Off-Slice Defects and Quadratic Macrostate Scale]\label{thm:xi-bulk-horizon-phi-area-law}
Let the golden ratio be \(\varphi=\frac{1+\sqrt5}{2}\), and define the \(\varphi\)-resolution ``horizon thickness scale''
\[
\epsilon_m^{\mathrm{hor}}:=\varphi^{-m}\qquad(m\ge 1).
\]
For any off-critical-line zero \(\rho=\beta+\mathrm{i}\gamma\) (\(0<\beta<\tfrac12\)), let \(\delta=\tfrac12-\beta>0\), and denote its endpoint layer remainder (horizon depth)
\(
h(\rho)=1-|w_\rho|^2
\)
as in Theorem \ref{thm:xi-offcritical-quadratic-radial-compression}.
Define the ``\(\varphi\)-resolution visibility threshold''
\[
m_\star(\rho):=\min\{m\in\NN:\ \epsilon_m^{\mathrm{hor}}\le h(\rho)\}.
\]
Then \(m_\star(\rho)\) has the closed form
\[
\boxed{\;
m_\star(\rho)=\left\lceil \log_{\varphi}\!\Bigl(\frac{\gamma^2+(1+\delta)^2}{4\delta}\Bigr)\right\rceil\; }.
\]
Moreover, the corresponding stable type scale (macrostate degrees of freedom) satisfies
\[
\card{X_{m_\star(\rho)}}\asymp \varphi^{m_\star(\rho)}
\asymp \frac{\gamma^2+(1+\delta)^2}{4\delta}
\sim \frac{\gamma^2}{4\delta}\qquad(|\gamma|\to\infty).
\]
\end{theorem}
\begin{proof}
By Theorem \ref{thm:xi-offcritical-quadratic-radial-compression},
\(
h(\rho)=\frac{4\delta}{\gamma^2+(1+\delta)^2}
\).
By definition, \(\epsilon_m^{\mathrm{hor}}\le h(\rho)\) is equivalent to
\(
\varphi^{-m}\le \frac{4\delta}{\gamma^2+(1+\delta)^2}
\),
i.e.,
\(
m\ge \log_{\varphi}\!\bigl(\frac{\gamma^2+(1+\delta)^2}{4\delta}\bigr)
\);
taking the smallest integer solution yields the stated \(m_\star(\rho)\).
From \(\card{X_m}=F_{m+2}\sim \varphi^m/\sqrt5\) (e.g., see the notational convention in Corollary \ref{cor:pom-maxfiber-rarity-two-measures} and the Fibonacci asymptotics used in its proof),
\(\card{X_{m_\star(\rho)}}\asymp \varphi^{m_\star(\rho)}\asymp 1/h(\rho)\), and substituting back the closed form yields the quadratic scaling law. \qed
\end{proof}

\begin{corollary}[Endpoint Fractal Index as the Power-Law Scale of ``Defect Visibility'']\label{cor:xi-bulk-horizon-defect-visibility-Dinfty}
Let the maximum fiber probability be \(p_{\max}(m)=\max_{x\in X_m}d_m(x)/2^m\) (Corollary \ref{cor:pom-max-fiber-rate-endpoint}), and let the R\'enyi dimension spectrum endpoint constant be
\[
D_\infty:=\lim_{q\to\infty}D_q=\frac{\log(2/\sqrt{\varphi})}{\log\varphi}
\]
as in Proposition \ref{prop:pom-rq-qinfty-endpoint}.
Then at the minimal visible resolution \(m_\star(\rho)\) of Theorem \ref{thm:xi-bulk-horizon-phi-area-law}, the following power-law scaling holds:
\[
\boxed{\;
\frac{1}{p_{\max}(m_\star(\rho))}
\asymp
\Bigl(\frac{1}{h(\rho)}\Bigr)^{D_\infty}
\asymp
\Bigl(\frac{\gamma^2+(1+\delta)^2}{4\delta}\Bigr)^{D_\infty}\; }.
\]
\end{corollary}
\begin{proof}
By Corollary \ref{cor:pom-max-fiber-rate-endpoint},
\(
\lim_{m\to\infty}-\frac{1}{m}\log p_{\max}(m)=\log(2/\sqrt{\varphi})
\),
so \(1/p_{\max}(m)\asymp (2/\sqrt{\varphi})^{m}\) (exponential scale).
By Theorem \ref{thm:xi-bulk-horizon-phi-area-law}, \(m_\star(\rho)=\lceil \log_\varphi(1/h(\rho))\rceil\); substituting yields
\[
\frac{1}{p_{\max}(m_\star(\rho))}
\asymp
\Bigl(\frac{2}{\sqrt{\varphi}}\Bigr)^{\log_{\varphi}(1/h(\rho))}
=\Bigl(\frac{1}{h(\rho)}\Bigr)^{\log(2/\sqrt{\varphi})/\log\varphi}
=\Bigl(\frac{1}{h(\rho)}\Bigr)^{D_\infty},
\]
and using \(1/h(\rho)=\frac{\gamma^2+(1+\delta)^2}{4\delta}\) gives the second equivalent expression. \qed
\end{proof}

\begin{proposition}[Ramanujan--Horizon Positivity--Toeplitz Certificate Transfer (Unified Finite-/Infinite-Dimensional Framework)]\label{prop:xi-ramanujan-horizon-toeplitz-transfer}
Under the regime of self-dual completion and unitary slice in this paper, RH-type zero constraints can be uniformly stated as ``horizon positivity certificates'':
\[
\text{(No disk-interior zeros)}
\ \Longleftrightarrow\
\text{(Carath\'eodory positivity)}
\ \Longleftrightarrow\
\text{(Toeplitz--PSD hierarchy)},
\]
where the strict equivalence for \(\Xi_{\zeta}\) is given by Theorem \ref{thm:app-horizon-toeplitz-lmi}, and the ``disk-interior zero exclusion'' with Jensen defect/repulsion radius/band exclusion single-parameter certificate form is closed by Theorem \ref{thm:xi-horizon-purity-criterion} and Theorem \ref{thm:xi-jensen-single-radius-band-exclusion}.
\par
In finite-dimensional rational/polynomial realizations, the same certificate language can be fully finitized: for example, the ``unit circle root'' condition for a self-inversive completion polynomial is equivalent to the vanishing of its Jensen defect (Theorem \ref{thm:xi-selfinversive-jensen-criterion}), and can be implemented via the Schur--Cohn/PSD interface as an auditable finite-level criterion (Remark \ref{rem:xi-schur-cohn-sdp-interface}).
In the graph-theoretic Ihara \(\zeta\) regular graph case, ``graph RH'' is equivalent to the Ramanujan spectral bound (see the surveys \cite{TerrasZetaGraphs,StarkTerras1996GraphCoverings}), constituting a classical external anchor for this finite-dimensional unified framework.
\par
Furthermore, in scattering theory on modular surfaces, the scattering coefficients (or scattering determinant) explicitly contain the ratio \(\zeta(2s-1)/\zeta(2s)\), thereby interpreting the zeros/poles of \(\zeta\) as analytic continuation singularities of scattering resonances (see, e.g., \cite{Iwaniec2002SpectralMethods}); this scattering anchor shows that ``Ramanujan-type unitary slice (boundary unitarity) \(\Rightarrow\) horizon positivity certificate'' is not restricted to graphs/finite kernels, but is naturally compatible with the \(\Xi\)/Toeplitz certificate language on infinite-dimensional analytic objects.
\end{proposition}

\paragraph{Alignment with Subsequent Subsections}
The above six blocks are concretized in subsequent subsections by the following numbered chains: block (1) on self-dual compression and phase closure appears in \S\ref{subsec:xi-from-self-dual}, \S\ref{subsubsec:ramanujan-phase-lift-symmetry}, and \S\ref{subsubsec:selfdual-blaschke-defect-fixed-slice};
block (2) on the horizon Carath\'eodory/Herglotz and Toeplitz--PSD hierarchy appears in Appendix \ref{subsubsec:app-horizon-weyl-carath-toeplitz};
block (3) on interval roots/Schur--Cohn and the Artin channel pipeline appears in the completion criteria of \S\ref{subsec:xi-from-self-dual}, Remark \ref{rem:xi-schur-cohn-sdp-interface}, and Remark \ref{rem:xi-artin-channel-pipeline};
block (4) on the \(d=6\) cut-and-project background and window-6 root system dictionary appears in \S\ref{subsec:icosahedral-6d-model-set} and \S\ref{subsec:bdry-tower-zeck-gut};
block (5) on R\'enyi/pressure/ledger and Chebotarev budget appears in Section \ref{sec:pom} and Theorem \ref{thm:pom-profinite-axis};
block (6) on endpoint compression, defect certificates, and the audit closed loop appears in \S\ref{subsec:dephysicalized-horizon-timefiber-nohair}, \S\ref{subsubsec:xi-ramanujan-horizon-purity-spectral-reduction}, and \S\ref{subsubsec:xi-fixed-slice-audit-chain},
while the falsifiability interfaces of the two unification conjectures appear in Section \ref{sec:discussion}.
